\chapter{Experimental Setup}
\label{chap:experiments}

This chapter summarizes the experimental setup used to evaluate both the re-identification (classification) pipeline and the population counting module.

\section{Datasets and split protocols}
Experiments are conducted primarily on datasets from the WildlifeReID-10k collection~\cite{adam2025wildlifereid10kwildlifereidentificationdataset} as described in Chapter~\ref{chap:dataset}. When using MegaDescriptor-L-384 for global embeddings, we follow the split protocol discussed in Chapter~\ref{chap:dataset} to avoid train/test leakage on datasets that were part of the MegaDescriptor training set.

\section{Evaluation metrics}
\subsection{Classification (re-identification)}
For the closed-set identification task, we report:
\begin{itemize}
    \item \textbf{Top-1 accuracy}: fraction of queries whose predicted identity matches the ground truth.
    \item \textbf{Top-5 accuracy}: fraction of queries where the correct identity appears among the five most confident predictions.
    \item \textbf{F1 score}: reported for completeness as a balance of precision and recall under the per-query identity decision rule.
\end{itemize}

\subsection{Population counting}
For HITL-NIS, we report:
\begin{itemize}
    \item \textbf{Estimated population size} \(\widehat{K}\) and its \textbf{standard error} (and the derived confidence interval).
    \item \textbf{Oracle query diagnostics}: number of oracle calls and unique queried pairs.
    \item \textbf{Runtime}: wall-clock time for the sampling-and-oracle loop (excluding one-time feature extraction and calibration training).
\end{itemize}

\section{Runtime measurement}
Runtime values reported in Chapter~\ref{chap:results} correspond to the wall-clock time measured by the evaluation scripts. For the classification pipeline, this includes dataset loading and matching, and may include cached feature reuse depending on whether features were previously computed. For population counting, the reported runtime focuses on the HITL-NIS loop itself, as this is the component that scales with the number of oracle queries.